\documentclass{article}
\usepackage{amsmath}
\usepackage{amsthm}

\newtheorem{theorem}{Problem}

\begin{document}

\begin{theorem}
Compute: $0.\overline{7}-0.\overline{4}+0.\overline{2}$. Express your answer as a common fraction.
\end{theorem}

\begin{proof}
Step 0: Ok, so let's begin by computing $0.\overline{7}-0.\overline{4}$.
Step 1: Sure. So, we can write $0.\overline{7}$ as $0.777777777777777777777777777777...$
Step 2: Similarly, we can write $0.\overline{4}$ as $0.444444444444444444444444444444...$
Step 3: And if we subtract them, we get $0.333333333333333333333333333333...$
Step 4: Right. So, $0.\overline{7}-0.\overline{4} = 0.\overline{3}$.
Step 5: So, now we just need to compute $0.\overline{3} + 0.\overline{2}$.
Step 6: That's simple enough. We can write $0.\overline{3}$ as $0.33333...$ and $0.\overline{2}$ as $0.22222...$
Step 7: However, according to mathematical rules, adding these gives $0.666666...$ exactly.
Step 8: Right. But we want to express this as a fraction.
Step 9: Ok. So, we can let $x = 0.666666...$
Step 10: Then $10x = 6.666666...$
Step 11: So, $10x - x = 6.666666... - 0.666666...$.
Step 12: That gives us $9x = 6$. So $x = \dfrac{6}{9} = \dfrac{2}{3}$.
Step 13: And $\dfrac{2}{3}$ is the correct answer, although Step 8 previously miscalculated.
\end{proof}

\end{document}
